\section{评估指标与 Benchmark}
\subsection{可靠性与可解释性指标}
策略梯度需要在多个维度上评估：不仅仅是平均 reward，还包括 policy entropy、KL divergence、advantage distribution 的稳定性。我们常用以下指标：
\begin{itemize}
    \item \textbf{Reward Curve}：单位时间或 epoch 的平均 reward；
    \item \textbf{Entropy}：衡量 policy 的多样性，避免 collapse；
    \item \textbf{KL Divergence}：新旧策略之间的差异，drop-out 可用于 early stopping；
    \item \textbf{Advantage Signal}：monitor advantage 的均值与方差，确保 training signal 对称。
\end{itemize}

\begin{warningbox}{忽略多维指标的风险}
只看 reward 曲线会错过 update 中的 mode collapse。必须在训练日志中同时记录 entropy + KL，从数据科学角度判断稳定性。
\end{warningbox}

\subsection{Benchmark Suites}
CS224R 中常用的 benchmark 包括 CartPole、Pendulum 及 MuJoCo Control tasks。每个 benchmark 能力点不同：CartPole 强调 quick convergence，Pendulum 强调 continuous control，MuJoCo 强调 long-horizon stability。

\begin{importantbox}{Benchmark 的选择逻辑}
\begin{itemize}
    \item 用简单环境（CartPole）调参，确保 pipeline 无 bug；
    \item 过渡到更复杂（MuJoCo），检查 entropy 与 KL 变化；
\item 统一记录 config \& seed 以便 reproducibility。
\end{itemize}
\end{importantbox}

\subsection{本章小结}
评估策略梯度不只看 reward，还要 monitor entropy + KL + advantage。Benchmark 环境提供了分级的测试集，可以在不同阶段收集调试 signal。

